\section{Introduction}

Weak-coupling lattice gauge fields strongly favor plaquette holonomies near
the identity.  For the heat-kernel action this statement is especially
transparent: the logarithmic cost of a plaquette $V$ is, to leading order,
its squared Riemannian distance from the identity.  It is tempting to promote
a small one-plaquette tail to a local statement that remains valid after
conditioning on every surrounding microscopic link configuration.  Such a
promotion would make dependence inexpensive: a boundary-uniform conditional
estimate could be inserted directly into a product-domination argument of
the type developed for dependent random fields
\cite{LiggettSchonmannStacey1997}, or into a Dobrushin interdependence
criterion \cite{Dobrushin1968}.

The quantifier ``every surrounding configuration'' is the issue.  A link in
a $d$-dimensional hypercubic lattice is shared by $m=2(d-1)$ plaquettes.  The
conditional link law minimizes the sum of all $m$ incident plaquette costs,
not each cost separately.  A central boundary defect can make these costs
incompatible: $m-1$ plaquettes request one link value while the remaining
plaquette requests a centrally displaced value.  Weak coupling then sharpens
the compromise instead of removing it.

This paper makes that observation into a finite-lattice theorem.  Let $G$ be
a compact connected Lie group, let $z\in\cZ(G)$ be nontrivial, and let $p_t$
be the heat kernel for $\tfrac12\Delta$ in a fixed bi-invariant metric.  On a
side-two open block we explicitly realize an incident staple pattern with
$m-1$ identity staples and one $z$-staple.  After conditioning on all other
links, the selected link $U$ has density
\begin{equation}\label{eq:intro-density}
  \gamma_t^z(\dd U)=Z_t(z)^{-1}p_t(U)^{m-1}p_t(z^{-1}U)\,\Haar.
\end{equation}
Varadhan's logarithmic heat-kernel asymptotic
\cite{Varadhan1967Fundamental,Varadhan1967Diffusion} converts
\eqref{eq:intro-density} into the energy
\begin{equation}\label{eq:intro-energy}
  F_z(U)=(m-1)\dist(\one,U)^2+\dist(z,U)^2.
\end{equation}
If $L=\dist(\one,z)$, the triangle inequality and a one-variable weighted
square identity imply
\[
 \min_U F_z(U)=\frac{m-1}{m}L^2,
 \qquad
 \dist(\one,U)=\frac Lm,
 \qquad
 \dist(z,U)=\frac{m-1}{m}L
\]
at every minimizer.  Multiple minimizing geodesics and cut-locus degeneracy
cause no difficulty: the set of minimizers may be large, but both displayed
distances are fixed on it.

Our first conclusion is an exact logarithmic asymptotic.  If the frustrated
plaquette is good when its distance from the identity is less than $\rho$ and
\begin{equation}\label{eq:intro-threshold}
  0<\rho<\frac{m-1}{m}L,
\end{equation}
then
\begin{equation}\label{eq:intro-rate}
 \lim_{t\downarrow0}2t\log\gamma_t^z
 \{\dist(z,U)<\rho\}
 =-m\left(\frac{m-1}{m}L-\rho\right)^2.
\end{equation}
Thus the conditional bad probability tends to one.  The construction is not
confined to a measure-zero boundary point: the rate gap persists uniformly
on an open neighborhood of the center-valued assignment, which has positive
product Haar measure.

Our second conclusion concerns the classical single-site Dobrushin matrix.
Each of the $m$ incident plaquettes has a distinct link opposite the selected
link.  Changing just that opposite link from $\one$ to a suitable central
value changes the coherent conditional law $p_t(U)^m\Haar$ into
\eqref{eq:intro-density}.  These two laws concentrate on disjoint sets, so
their total variation distance tends to one.  Repeating the comparison for
the $m$ distinct opposite links makes the selected Dobrushin row sum tend to
at least $m$; hence the usual row-sum criterion is violated for all
sufficiently small $t$.  This statement concerns the criterion, not Gibbs
uniqueness itself.

For $G=\SUtwo$, $z=-\one$, and the conventional metric normalization
$L=\pi$.  In four dimensions $m=6$ and the conditional class-angle energy is
\[
  5\theta^2+(\pi-\theta)^2
  =6\left(\theta-\frac\pi6\right)^2+\frac{5\pi^2}{6}.
\]
The selected link concentrates at class angle $\pi/6$, while the frustrated
plaquette concentrates at $5\pi/6$.  This recovers the concrete case that
motivates the general result, but neither the proof nor the theorem depends
on numerical quadrature.

The heat-kernel lattice action was introduced as a gauge-theory action by
Menotti and Onofri \cite{MenottiOnofri1981}.  Central twists and vortex free
energies have a separate, extensive history
\cite{tHooft1979,MackPetkova1980,deForcrandVonSmekal2002}; those works study
global flux sectors and their free-energy costs.  Our theorem is local and
conditional: it identifies what happens when the cost of a frustrating
boundary is erased by a supremum over microscopic boundary data.  It does
not assert that center frustration is new, nor does it replace the
small-field/large-field decompositions in constructive renormalization
\cite{Balaban1987,Balaban1988,Balaban1989}.

Section~\ref{sec:setup} fixes lattice, heat-kernel, and conditional-probability
conventions.  Section~\ref{sec:main-results} proves the central realization,
the conditional large-deviation theorem, and the exact rate.
Section~\ref{sec:stability-dobrushin} establishes positive-measure stability
and the Dobrushin consequence.  Section~\ref{sec:extensions} treats multiple
frustrated staples, low dimensions, and the $\SUtwo$ specialization.
Sections~\ref{sec:related} and \ref{sec:scope} locate the result and delimit
its implications.  The appendices supply the conditional-kernel derivation,
uniform Laplace bounds, incidence certificate, stress tests, and
reproducibility information.
